% !TEX root = ../main.tex
\lecture{3}{2026-10-01}{Paretto und Präferenzen}

\section{Modellökonomien}

\begin{definition}
    [Paretto Verbesserung] änderungen in der Warenverteilung/Resourcenallokation, die keine der Mitglieder benachteiligen und mindestens einen Mitglieder eine bessere Position erlauben relativ zu einer gegebnen Situation (\textit{Status Quo })\footnote{Ist Situation Zero-Sum game oder Plus-Sum}
\end{definition}

\begin{definition}
    [Ineffizienz] wenn die Möglichkeit einer Verbesserung der Verteilung von Resourcen an Bedürfnissen möglich ist, jedoch nicht durchgeführt wird.
\end{definition}

\subsection{Zulässige Allokationen}
Bei \textit{Knappheit} welche Menge möglicher allokationen entsteht. Innerhalb der Menge möglicher Allokationen, sind nur begrenzte \textit{Verbesserungen} möglich, bis die materielle Grenze erreicht ist. An der Grenze ist \textit{\textbf{keine}} Parettoverbesserung möglich.

\subsubsection{Wie wurde Status Quo der Verteilung erreicht?}
Bei einer grenzwertigen Verteilung ist keine Paretto verbesserung möglich, jemand muss verlieren.

\subsection{Präferenzen}

\begin{definition}
    [Präferenz] individuelle Vorliebe. Es besteht eine Rangordnung von Präferenzen basierend auf verschiedene vergleichbare Kriteren. \footnote{Ausgehend vom \textit{lebendigen Ganzen}, das konkret zu sein scheint, werden einfache Kategorien, Konzepte und Regeln hergeleitet zu einem System}
\end{definition}

\subsubsection{Güterbündel}
Bei gegebnem Wahl zweier oder mehrere Güterbündel ($X_n$), kann formell eine Rangordnung defnitiert werden.

\textbf{Bevorzugen}
\begin{equation}
    X_1 \succ X_2
\end{equation}

\textbf{schwach Bevorzugen oder gleich}
\begin{equation}
    X_1 \succeq X_2
\end{equation}

\textbf{Indifferenz}
\begin{equation}
    X_1 \sim X_2
\end{equation}

\footnote{Es kann nicht quantifiziert werden, nur Kategorisiert. Güter sind nicht reduzierbar auf gleiche und quantifizierbare Grössen}

\textbf{Annahmen}
\begin{itemize}
    \item Vollständigkeit
    \item Transitivität
    \item Stetigkeit
    \item Nichtsättigung
    \item Konvexität
\end{itemize}


\subsection{Angewandte Mikroökonomie}

Empirische Forschung und beobachtung. \footnote{Die Empirie ist aber...}

\subsubsection{Matching-Market}
\textbf{Beispiele}
\begin{itemize}
    \item Partnersuche
    \item Organspenden
\end{itemize}

\textbf{Eigenschaften}
\begin{itemize}
    \item Keine Preisbildung
    \item Rangordnung von Präferenzen werden gebraucht
\end{itemize}

\subsubsection{Partnersuche}

\textbf{Fragen}
\begin{itemize}
    \item Welche Präferenz zeigen sich bei Männer und Frauen
    \item Wie verändern sich diese in Relation zum Angebot (Potenzielle Partner)
    \item Zufällige Veränderung des Partnerpools, sollten Zufällige ver
\end{itemize}

\section{Konsument und Nachfrage}


\subsection{Budgetbeschränkung}

\subsubsection{Voraussetzungen}
\begin{itemize}
    \item montliches Einkommen
    \item Zwei Waren Wohnraum W ($m^2$ / Monat) und Fertiggerichte F kg/monat
    \item Preise für Wohnraum $P_w$ und $P_f$
    \item \textit{Preisnehmerverhalten} keinen individuellen Einfluss auf Preise
    \item Güterbündel Tüpel $X_n = (W,F)$ \footnote{Alle wichtige Fragen scheinen im Hintergrund zu sein: Warum beträgt Einkommen X so viel, wieso kostet wohnraum x mal mehr als F, was ist ihnen überhaupt Gemeinsam für einen Quantitativen Vergleich }
\end{itemize}

\subsubsection{Budgetgerade}
\begin{definition}
    [Budgetbedingung] Was eine Person sich leisten \textit{kann}. (W,F) kostet $P_wW + P_fF$
    \begin{equation}
         P_wW+P_fF \le M
    \end{equation} wobei M = Budget
\end{definition}

\begin{definition}
    [Budgetgerade]
 \textbf{Bei voller Ausgabe vom Budget}

    \begin{equation}
        P_wW + P_fF = M
    \end{equation}
    \textbf{Umgeformt Budgetgerade}
    \begin{equation}
        F = \frac{M}{P_f} - \frac{P_w}{P_f}W
    \end{equation}
\end{definition}
Stellt alle Konvexkombinationen von W und F dar, worin Preis von F und Preis von W zusammen M ergeben. Randpunkte der Menge, wo F für W ZeroSum ausgetauscht werden müssen um im Budget zu passen. Randpunkte nicht unbedingt optimal, weil \textbf{Präfenzen} noch abstrahiert.

\subsubsection{Einseitige Preisänderungen}
Veränderung der Budget gerade bei Verändernden Preisen von Güter im Bündel.
\begin{itemize}
    \item Bei erhöhung von $P_wW$ und gleichbleibendem M, reduziert sich die Anzahl möglicher Konvexkombinationen
    \item Gleich bei $P_fF$
    \item Bei einer reduktion von M, wirkt sich auf Kurve so aus, als wären $P_f$ und $P_w$ gestiegen. Allgemeine reduktion der möglichen Kombinationen, steigerung der Punkte ausserhalb der Menge.
\end{itemize}

\subsubsection{Simultane Änderungen}
Änderungen können, sich gegenseitig neutralisieren. So: 
\begin{itemize}
    \item Verdopplung der Preise gleich Halbierung des Budgets
    \item So bei Verdopplung des Budgets und Verdopplung der Preise $\implies$ Gleiche Mengemächtigkeit vor und nach.
\end{itemize}

\subsection{Präferenzen und Nutzenfunktion }
\begin{definition}
    [Nutzenfunktion U] von den Tupel innerhalb der leistbaren Menge. Welcher ist optimal, gegeben Präferenzen und Rationalität.
    \begin{equation}
        U: \mathcal{X} \rightarrow \mathbb{R} \quad U(W,F) = WF
    \end{equation}
    Gegeben Güterbündel ($\mathcal{X}$) wird eine Reele Zahl zugeordnet, die den Nutzen messen soll.
    Interpretation ist \textit{\textbf{Ordinal}} also qualitativ, ob grösser oder kleiner Nutzen, \textbf{nicht} quantitativ $X_1 = 2 X_2$ \footnote{Also einfach ignorieren, dass die Resultate quantitativ vergleichbar sind als Ausgang von dem Widerspruch, dass Nützlichkeit verschiedener Dinge sich nicht quantifizierbar vergleichen lässt?}
\end{definition}

\subsubsection{Höhenlinien}
Nutzen ist eine multidimensionale Funktion. Wobei $n$ gleich der Anzahl Güter im Bündel ist. Auf einer gleichen Linie an der oberfläche der Nutzen ebene, haben Güterkombinationen denselben Nutzenniveau.
\begin{equation}
    U: \mathbb{R} \rightarrow \mathbb{R}^{n}
\end{equation}
